% ch2-linear-maps.tex — Chapter 2: Linear Maps and Matrices
% Definitions, kernel/image, rank-nullity theorem, matrix representations,
% change of basis, operations on matrices, dual spaces, applications.

\chapter{Linear Maps and Matrices}\label{ch:linear-maps}

% ============================================================
% Motivating introduction
% ============================================================

The study of vector spaces becomes truly powerful only when we understand
the \emph{maps} between them that respect their algebraic structure.
Consider the Euclidean plane~$\R^2$. Many natural geometric
transformations---rotations, reflections, projections onto a line,
scalings---share a remarkable property: they preserve the operations of
vector addition and scalar multiplication. A rotation through angle
$\theta$ sends the sum of two vectors to the sum of their rotated
images; scaling by a factor~$\lambda$ commutes with addition.

\begin{center}
\begin{tikzpicture}[>=Stealth, scale=1.2]
  % Original vectors
  \draw[->, thick, blue!70!black] (0,0) -- (2,0.5) node[right] {$\vec{v}$};
  \draw[->, thick, red!70!black] (0,0) -- (0.8,1.8) node[above] {$\vec{w}$};
  \draw[->, thick, black!60] (0,0) -- (2.8,2.3) node[above right] {$\vec{v}+\vec{w}$};

  % Rotation arrow
  \draw[->, thick, dashed, gray] (3.5,1.2) to[bend left=20]
    node[midway,above] {\small rotation $R_\theta$} (5.5,1.2);

  % Rotated vectors (approx 45 degrees)
  \begin{scope}[shift={(7,0)}, rotate=40]
    \draw[->, thick, blue!70!black] (0,0) -- (2,0.5)
      node[right] {$R_\theta(\vec{v})$};
    \draw[->, thick, red!70!black] (0,0) -- (0.8,1.8)
      node[above] {$R_\theta(\vec{w})$};
    \draw[->, thick, black!60] (0,0) -- (2.8,2.3)
      node[above right] {$R_\theta(\vec{v}+\vec{w})$};
  \end{scope}
\end{tikzpicture}
\captionof{figure}{A rotation maps the parallelogram law to a rotated
  parallelogram: the image of a sum is the sum of the images.}
\label{fig:rotation-intro}
\end{center}

These structure-preserving maps are called \emph{linear maps}. They form
the bridge between the abstract theory of vector spaces and the concrete
computations of matrix algebra. In this chapter we develop their theory
systematically: we define linear maps, study their kernel and image,
prove the fundamental rank--nullity theorem, represent linear maps by
matrices, and learn to change bases. Along the way we shall see that
matrix multiplication is nothing but the composition of linear maps in
disguise.

Throughout this chapter, $\K$ denotes a field (typically $\R$ or $\C$),
and $E$, $F$, $G$ denote $\K$-vector spaces unless stated otherwise.
We freely use results from 
ef{ch:vector-spaces}, especially the
notions of subspace, basis, and dimension.


% ============================================================
\section{Definition and first properties}\label{sec:linear-map-def}
% ============================================================

\begin{definition}{Linear map}{def:linear-map}
\index{linear map}
Let $E$ and $F$ be vector spaces over~$\K$. A map
$f\colon E \to F$ is \emph{linear} (or is a \emph{$\K$-linear map},
or a \emph{homomorphism of vector spaces}) if it satisfies the
following two conditions:
\begin{enumerate}[label=(\roman*)]
  \item \textbf{Additivity:}
    $f(u+v) = f(u) + f(v)$\quad for all $u,v\in E$.
  \item \textbf{Homogeneity:}
    $f(\lambda u) = \lambda\, f(u)$\quad for all $\lambda\in\K$,
    $u\in E$.
\end{enumerate}
Equivalently, $f$ is linear if and only if
$f(\lambda u + \mu v) = \lambda\, f(u) + \mu\, f(v)$
for all $\lambda,\mu\in\K$ and all $u,v\in E$.
\end{definition}

\begin{remark}{Terminology}{rem:terminology}
\index{endomorphism}\index{isomorphism}\index{automorphism}
Several special names are used:
\begin{itemize}
  \item A linear map $f\colon E\to E$ (same domain and codomain) is
    called an \emph{endomorphism} of~$E$.
  \item A bijective linear map is called an \emph{isomorphism}.
  \item A bijective endomorphism is called an \emph{automorphism}.
  \item A linear map $f\colon E\to\K$ (codomain is the base field)
    is called a \emph{linear form}\index{linear form}.
\end{itemize}
\end{remark}

\begin{proposition}{Elementary properties}{prop:linear-elementary}
Let $f\colon E\to F$ be a linear map. Then:
\begin{enumerate}[label=(\roman*)]
  \item $f(\vzero_E) = \vzero_F$.
  \item $f(-v) = -f(v)$ for all $v\in E$.
  \item $f\!\left(\sum_{i=1}^{n} \lambda_i v_i\right)
    = \sum_{i=1}^{n} \lambda_i\, f(v_i)$
    for all $\lambda_i\in\K$, $v_i\in E$.
\end{enumerate}
\end{proposition}

\begin{proof}
(i) Setting $\lambda = 0$ in homogeneity: $f(\vzero_E) = f(0\cdot v) =
0\cdot f(v) = \vzero_F$.

(ii) $f(-v) = f((-1)\cdot v) = (-1)\cdot f(v) = -f(v)$.

(iii) By induction on~$n$, using additivity and homogeneity at each step.
The base case $n=1$ is homogeneity, and the inductive step follows from
$f\!\left(\sum_{i=1}^{n+1}\lambda_i v_i\right)
= f\!\left(\sum_{i=1}^{n}\lambda_i v_i + \lambda_{n+1}v_{n+1}\right)
= f\!\left(\sum_{i=1}^{n}\lambda_i v_i\right) + \lambda_{n+1}\,f(v_{n+1})$.
\end{proof}

\begin{proposition}{A linear map is determined by its values on a basis}%
{prop:linear-determined-by-basis}
\index{linear map!determined by basis}
Let $\cB = (\vece{1},\ldots,\vece{n})$ be a basis of~$E$, and let
$w_1,\ldots,w_n$ be arbitrary vectors in~$F$. Then there exists a
\emph{unique} linear map $f\colon E\to F$ such that
$f(\vece{i}) = w_i$ for all $i\in\set{1,\ldots,n}$.
\end{proposition}

\begin{proof}
\textit{Existence.}\; Every $v\in E$ has a unique decomposition
$v = \sum_{i=1}^{n}\lambda_i\vece{i}$. Define
$f(v) \deff \sum_{i=1}^{n}\lambda_i w_i$.
One verifies directly that $f$ is linear: for
$v = \sum \lambda_i \vece{i}$ and $v' = \sum \mu_i \vece{i}$ and
$\alpha,\beta\in\K$,
\[
  f(\alpha v + \beta v')
  = f\!\Bigl(\sum_{i}(\alpha\lambda_i+\beta\mu_i)\vece{i}\Bigr)
  = \sum_{i}(\alpha\lambda_i+\beta\mu_i)w_i
  = \alpha\sum_{i}\lambda_i w_i + \beta\sum_{i}\mu_i w_i
  = \alpha f(v)+\beta f(v').
\]
Moreover $f(\vece{i}) = w_i$ since
$\vece{i} = 0\cdot\vece{1}+\cdots+1\cdot\vece{i}+\cdots+0\cdot\vece{n}$.

\textit{Uniqueness.}\; Suppose $g\colon E\to F$ is linear with
$g(\vece{i}) = w_i$. For any $v = \sum \lambda_i\vece{i}$,
$g(v) = \sum \lambda_i\, g(\vece{i}) = \sum \lambda_i w_i = f(v)$.
\end{proof}


% ============================================================
\section{Examples of linear maps}\label{sec:linear-map-examples}
% ============================================================

\begin{example}{Rotation in $\R^2$}{ex:rotation}
\index{rotation}
For a fixed angle $\theta$, the map $R_\theta\colon\R^2\to\R^2$
defined by
\[
  R_\theta\!\begin{pmatrix}x\\y\end{pmatrix}
  = \begin{pmatrix}\cos\theta & -\sin\theta\\
    \sin\theta & \cos\theta\end{pmatrix}
    \begin{pmatrix}x\\y\end{pmatrix}
  = \begin{pmatrix}x\cos\theta - y\sin\theta\\
    x\sin\theta + y\cos\theta\end{pmatrix}
\]
is a linear map (an automorphism, in fact, with inverse~$R_{-\theta}$).
\end{example}

\begin{example}{Orthogonal projection onto a line}{ex:projection-line}
\index{projection!orthogonal}
Let $\ell\subset\R^2$ be the line through the origin with unit
direction vector $\vec{u} = (\cos\alpha,\sin\alpha)^{\mathsf{T}}$.
The orthogonal projection onto $\ell$ is
\[
  \proj_\ell(\vec{v}) = \inner{\vec{v},\vec{u}}\,\vec{u}
  = \begin{pmatrix}\cos^2\alpha & \cos\alpha\sin\alpha\\
    \cos\alpha\sin\alpha & \sin^2\alpha\end{pmatrix}
  \begin{pmatrix}x\\y\end{pmatrix}.
\]
This is linear. Note $\proj_\ell \circ \proj_\ell = \proj_\ell$
(idempotent).
\end{example}

\begin{example}{Reflection}{ex:reflection}
\index{reflection}
The reflection of $\R^2$ across the line $\ell$ through the origin
with direction angle~$\alpha$ is given by
\[
  S_\alpha\!\begin{pmatrix}x\\y\end{pmatrix}
  = \begin{pmatrix}\cos 2\alpha & \sin 2\alpha\\
    \sin 2\alpha & -\cos 2\alpha\end{pmatrix}
  \begin{pmatrix}x\\y\end{pmatrix}.
\]
This is a linear map satisfying $S_\alpha^2 = \Id$.
\end{example}

\begin{example}{Scaling (homothety)}{ex:scaling}
\index{scaling}\index{homothety}
For $\lambda\in\K$, the map $h_\lambda\colon E\to E$ defined by
$h_\lambda(v) = \lambda v$ is linear. It contracts ($|\lambda|<1$),
dilates ($|\lambda|>1$), or reflects ($\lambda<0$) vectors.
\end{example}

\begin{example}{Differentiation}{ex:differentiation}
\index{differentiation}
Let $\cP_n(\R)$ be the space of real polynomials of degree at most~$n$.
The derivative map $D\colon\cP_n(\R)\to\cP_{n-1}(\R)$ defined by
$D(p) = p'$ is linear, since $(p+q)' = p'+q'$ and $(\lambda p)' =
\lambda p'$.
\end{example}

\begin{example}{Integration}{ex:integration}
\index{integration}
The map $I\colon\mathcal{C}([0,1],\R)\to\R$ defined by
$I(f) = \int_0^1 f(t)\,dt$ is a linear form on the vector space of
continuous real-valued functions on $[0,1]$.
\end{example}

\begin{example}{Matrix--vector multiplication}{ex:matrix-vector}
\index{matrix!multiplication}
Let $A\in\Mat_{m,n}(\K)$. The map $f_A\colon\K^n\to\K^m$ defined by
$f_A(x) = Ax$ is linear. Conversely, every linear map
$\K^n\to\K^m$ arises this way (see 
ef{sec:matrix-of-linear-map}).
\end{example}


% ============================================================
\section{Geometric illustrations}\label{sec:geometric-illustrations}
% ============================================================

The following figures illustrate how several linear maps transform
the unit square in~$\R^2$.

\begin{figure}[ht]
\centering
\begin{tikzpicture}[>=Stealth, scale=1.3]
  % --- Rotation ---
  \begin{scope}
    \fill[blue!10] (0,0) -- (1,0) -- (1,1) -- (0,1) -- cycle;
    \draw[blue!50, thick] (0,0) -- (1,0) -- (1,1) -- (0,1) -- cycle;
    \node[below] at (0.5,-0.1) {\small original};
  \end{scope}
  \begin{scope}[shift={(3,0)}, rotate=30]
    \fill[red!15] (0,0) -- (1,0) -- (1,1) -- (0,1) -- cycle;
    \draw[red!60, thick] (0,0) -- (1,0) -- (1,1) -- (0,1) -- cycle;
  \end{scope}
  \node[below] at (3.3,-0.2) {\small rotation ($30^\circ$)};

  % --- Projection ---
  \begin{scope}[shift={(6.5,0)}]
    \fill[blue!10] (0,0) -- (1,0) -- (1,1) -- (0,1) -- cycle;
    \draw[blue!50, thick] (0,0) -- (1,0) -- (1,1) -- (0,1) -- cycle;
    \node[below] at (0.5,-0.1) {\small original};
  \end{scope}
  \begin{scope}[shift={(9.5,0)}]
    \draw[green!60!black, very thick] (0,0) -- (1,0);
    \draw[green!60!black, very thick] (0,0) -- (1,0)
      node[midway, below=4pt] {\small proj.\ onto $x$-axis};
    \fill[green!20, opacity=0.5] (0,0) -- (1,0) -- (1,0) -- (0,0);
  \end{scope}
\end{tikzpicture}

\bigskip

\begin{tikzpicture}[>=Stealth, scale=1.3]
  % --- Shear ---
  \begin{scope}
    \fill[blue!10] (0,0) -- (1,0) -- (1,1) -- (0,1) -- cycle;
    \draw[blue!50, thick] (0,0) -- (1,0) -- (1,1) -- (0,1) -- cycle;
    \node[below] at (0.5,-0.1) {\small original};
  \end{scope}
  \begin{scope}[shift={(3,0)}]
    \fill[orange!15] (0,0) -- (1,0) -- (1.5,1) -- (0.5,1) -- cycle;
    \draw[orange!70, thick] (0,0) -- (1,0) -- (1.5,1) -- (0.5,1) -- cycle;
    \node[below] at (0.75,-0.1) {\small shear ($k=0.5$)};
  \end{scope}

  % --- Scaling ---
  \begin{scope}[shift={(6.5,0)}]
    \fill[blue!10] (0,0) -- (1,0) -- (1,1) -- (0,1) -- cycle;
    \draw[blue!50, thick] (0,0) -- (1,0) -- (1,1) -- (0,1) -- cycle;
    \node[below] at (0.5,-0.1) {\small original};
  \end{scope}
  \begin{scope}[shift={(9.5,0)}, scale=0.6]
    \fill[violet!15] (0,0) -- (1,0) -- (1,1) -- (0,1) -- cycle;
    \draw[violet!60, thick] (0,0) -- (1,0) -- (1,1) -- (0,1) -- cycle;
  \end{scope}
  \node[below] at (9.8,-0.1) {\small scaling ($\lambda=0.6$)};
\end{tikzpicture}

\caption{Effect of linear maps on the unit square: rotation, projection
  onto the $x$-axis, horizontal shear, and uniform scaling.}
\label{fig:linear-map-effects}
\end{figure}


% ============================================================
\section{Kernel and image}\label{sec:kernel-image}
% ============================================================

\begin{definition}{Kernel}{def:kernel}
\index{kernel}
Let $f\colon E\to F$ be a linear map. The \emph{kernel} (or
\emph{null space}) of~$f$ is
\[
  \Ker f \deff \setcond{v\in E}{f(v) = \vzero_F}.
\]
\end{definition}

\begin{definition}{Image}{def:image-linear}
\index{image!of a linear map}
The \emph{image} (or \emph{range}) of~$f$ is
\[
  \im f \deff \setcond{f(v)}{v\in E}
  = \setcond{w\in F}{\exists\, v\in E,\; f(v) = w}.
\]
\end{definition}

\begin{theorem}{Kernel and image are subspaces}{thm:ker-im-subspaces}
Let $f\colon E\to F$ be linear. Then:
\begin{enumerate}[label=(\roman*)]
  \item $\Ker f$ is a subspace of~$E$.
  \item $\im f$ is a subspace of~$F$.
\end{enumerate}
\end{theorem}

\begin{proof}
(i) We verify the three subspace criteria.
\emph{Non-empty:} $f(\vzero_E) = \vzero_F$, so
$\vzero_E\in\Ker f$.
\emph{Closed under addition:} if $u,v\in\Ker f$, then
$f(u+v) = f(u)+f(v) = \vzero_F+\vzero_F = \vzero_F$, so $u+v\in\Ker f$.
\emph{Closed under scalar multiplication:} if $v\in\Ker f$ and
$\lambda\in\K$, then
$f(\lambda v) = \lambda f(v) = \lambda\cdot\vzero_F = \vzero_F$,
so $\lambda v\in\Ker f$.

(ii) \emph{Non-empty:} $\vzero_F = f(\vzero_E)\in\im f$.
\emph{Closed under addition:} if $w_1,w_2\in\im f$, write
$w_i = f(v_i)$; then $w_1+w_2 = f(v_1)+f(v_2) = f(v_1+v_2)\in\im f$.
\emph{Scalar multiplication:}
$\lambda w_1 = \lambda f(v_1) = f(\lambda v_1)\in\im f$.
\end{proof}

\begin{figure}[ht]
\centering
\begin{tikzpicture}[>=Stealth, scale=1.0]
  % E ellipse
  \draw[thick, blue!60] (0,0) ellipse (2cm and 3cm);
  \node[above, blue!80!black, font=\large] at (0,3.2) {$E$};

  % Ker f inside E
  \fill[red!15, opacity=0.7] (0,-0.5) ellipse (1cm and 1.5cm);
  \draw[thick, red!60] (0,-0.5) ellipse (1cm and 1.5cm);
  \node[red!70!black, font=\small] at (0,-0.5) {$\Ker f$};

  % F ellipse
  \draw[thick, green!50!black] (7,0) ellipse (2cm and 3cm);
  \node[above, green!60!black, font=\large] at (7,3.2) {$F$};

  % Im f inside F
  \fill[orange!20, opacity=0.7] (7,0.3) ellipse (1.2cm and 1.8cm);
  \draw[thick, orange!70] (7,0.3) ellipse (1.2cm and 1.8cm);
  \node[orange!70!black, font=\small] at (7,0.3) {$\im f$};

  % zero in F
  \fill (7,-1.2) circle (2pt) node[right, font=\small] {$\vzero_F$};

  % Arrow from Ker f to 0_F
  \draw[->, thick, red!60, dashed] (1,-0.5) to[bend left=15]
    node[midway, above, font=\footnotesize] {$f$} (6.3,-1.2);

  % Arrow from E to Im f
  \draw[->, thick, blue!60] (1.5,1.5) to[bend left=20]
    node[midway, above, font=\footnotesize] {$f$} (5.8,1.5);
\end{tikzpicture}
\caption{The kernel is a subspace of~$E$ mapping to~$\vzero_F$; the
  image is a subspace of~$F$.}
\label{fig:kernel-image}
\end{figure}

\begin{proposition}{Image of a Spanning set}{prop:image-Spanning}
If $E = \Span(v_1,\ldots,v_p)$, then
$\im f = \Span(f(v_1),\ldots,f(v_p))$.
In particular, if $\cB = (\vece{1},\ldots,\vece{n})$ is a basis
of~$E$, then $\im f = \Span(f(\vece{1}),\ldots,f(\vece{n}))$.
\end{proposition}

\begin{proof}
Every $w\in\im f$ has the form $w = f(v)$ where
$v = \sum_{i=1}^{p}\lambda_i v_i$. By linearity,
$w = \sum \lambda_i f(v_i) \in \Span(f(v_1),\ldots,f(v_p))$.
Conversely, each $f(v_i)\in\im f$, and $\im f$ is a subspace, so it
contains $\Span(f(v_1),\ldots,f(v_p))$.
\end{proof}

\begin{definition}{Rank and nullity}{def:rank-nullity}
\index{rank}\index{nullity}
Let $f\colon E\to F$ be a linear map between finite-dimensional spaces.
\begin{itemize}
  \item The \emph{rank} of $f$ is $\rank(f) \deff \dim(\im f)$.
  \item The \emph{nullity} of $f$ is $\dim(\Ker f)$.
\end{itemize}
\end{definition}


% ============================================================
\section{The rank--nullity theorem}\label{sec:rank-nullity}
% ============================================================

\begin{theorem}{Rank--nullity theorem (dimension theorem)}%
{thm:rank-nullity}
\index{rank--nullity theorem}\index{dimension theorem}
Let $E$ be a finite-dimensional vector space and let $f\colon E\to F$
be a linear map. Then $\im f$ is finite-dimensional and
\[
  \boxed{\dim E \;=\; \dim(\Ker f) \;+\; \dim(\im f)
    \;=\; \dim(\Ker f) \;+\; \rank(f).}
\]
\end{theorem}

\begin{proof}
Set $n = \dim E$ and $r = \dim(\Ker f)$.
Let $(u_1,\ldots,u_r)$ be a basis of $\Ker f$. Since $\Ker f$ is a
subspace of the finite-dimensional space~$E$, we may extend this to a
basis of~$E$:
\[
  (u_1,\ldots,u_r,\, v_1,\ldots,v_s)
\]
where $s = n-r$. We claim that $(f(v_1),\ldots,f(v_s))$ is a basis
of $\im f$, which will prove $\dim(\im f) = s = n - r$.

\medskip
\textit{Step 1: $(f(v_1),\ldots,f(v_s))$ Spans $\im f$.}\;
Let $w\in\im f$, say $w = f(x)$ with
$x = \alpha_1 u_1 + \cdots + \alpha_r u_r +
\beta_1 v_1 + \cdots + \beta_s v_s$.
Then
\[
  w = f(x) = \alpha_1 \underbrace{f(u_1)}_{=\,\vzero}
  + \cdots + \alpha_r \underbrace{f(u_r)}_{=\,\vzero}
  + \beta_1 f(v_1) + \cdots + \beta_s f(v_s)
  = \beta_1 f(v_1) + \cdots + \beta_s f(v_s),
\]
so $w\in\Span(f(v_1),\ldots,f(v_s))$.

\medskip
\textit{Step 2: $(f(v_1),\ldots,f(v_s))$ is linearly independent.}\;
Suppose $\beta_1 f(v_1) + \cdots + \beta_s f(v_s) = \vzero_F$.
By linearity, $f(\beta_1 v_1 + \cdots + \beta_s v_s) = \vzero_F$,
so $\beta_1 v_1 + \cdots + \beta_s v_s \in \Ker f$.
Since $(u_1,\ldots,u_r)$ is a basis of $\Ker f$, there exist
$\gamma_1,\ldots,\gamma_r\in\K$ such that
\[
  \beta_1 v_1 + \cdots + \beta_s v_s
  = \gamma_1 u_1 + \cdots + \gamma_r u_r.
\]
Rearranging:
$\gamma_1 u_1 + \cdots + \gamma_r u_r
- \beta_1 v_1 - \cdots - \beta_s v_s = \vzero_E$.
Since $(u_1,\ldots,u_r,v_1,\ldots,v_s)$ is a basis of~$E$, all
coefficients are zero. In particular, $\beta_1 = \cdots = \beta_s = 0$.

\medskip
We conclude that $(f(v_1),\ldots,f(v_s))$ is a basis of $\im f$,
so $\dim(\im f) = s = n - r$, i.e.,
$\dim E = \dim(\Ker f) + \dim(\im f)$.
\end{proof}

\begin{corollary}{Dimension bound on rank}{cor:rank-bound}
For any linear map $f\colon E\to F$ with $E$ finite-dimensional,
\[
  \rank(f) \leq \min\!\bigl(\dim E,\, \dim F\bigr).
\]
\end{corollary}

\begin{proof}
$\rank(f) = \dim(\im f) \leq \dim F$ since $\im f\subseteq F$, and
$\rank(f) = \dim E - \dim(\Ker f) \leq \dim E$.
\end{proof}


% ============================================================
\section{Injectivity, surjectivity, bijectivity}%
\label{sec:inject-surject-biject}
% ============================================================

\begin{theorem}{Injectivity criterion}{thm:inject-kernel}
\index{injectivity!of linear maps}
A linear map $f\colon E\to F$ is injective if and only if
$\Ker f = \set{\vzero_E}$.
\end{theorem}

\begin{proof}
$(\Rightarrow)$\; Suppose $f$ is injective. If $v\in\Ker f$, then
$f(v) = \vzero_F = f(\vzero_E)$. Injectivity gives $v = \vzero_E$.

$(\Leftarrow)$\; Suppose $\Ker f = \set{\vzero_E}$. If
$f(u) = f(v)$, then $f(u-v) = f(u)-f(v) = \vzero_F$, so
$u - v\in\Ker f = \set{\vzero_E}$, hence $u = v$.
\end{proof}

\begin{proposition}{Characterizations in finite dimension}%
{prop:iso-finite-dim}
Let $f\colon E\to F$ be linear with $\dim E = \dim F = n < \infty$.
Then the following are equivalent:
\begin{enumerate}[label=(\roman*)]
  \item $f$ is injective.
  \item $f$ is surjective.
  \item $f$ is bijective (i.e., an isomorphism).
  \item $\rank(f) = n$.
\end{enumerate}
\end{proposition}

\begin{proof}
By the rank--nullity theorem, $\rank(f) = n - \dim(\Ker f)$.

(i) $\Leftrightarrow$ (iv):
$f$ is injective $\Leftrightarrow$ $\Ker f = \set{\vzero}$
$\Leftrightarrow$ $\dim(\Ker f) = 0$ $\Leftrightarrow$ $\rank(f) = n$.

(ii) $\Leftrightarrow$ (iv):
$f$ is surjective $\Leftrightarrow$ $\im f = F$
$\Leftrightarrow$ $\dim(\im f) = \dim F = n$ $\Leftrightarrow$
$\rank(f) = n$.

(iii) $\Leftrightarrow$ (i) $\wedge$ (ii): immediate from the
equivalences above.
\end{proof}

\begin{remark}{Same dimension is essential}{rem:same-dim}
The equivalence of injectivity and surjectivity \emph{fails} when
$\dim E\neq\dim F$. For instance, the inclusion
$\R^2\hookrightarrow\R^3$ (sending $(x,y)\mapsto(x,y,0)$) is
injective but not surjective.
\end{remark}

\begin{definition}{Isomorphism of vector spaces}{def:isomorphism}
\index{isomorphism}
Two vector spaces $E$ and $F$ over $\K$ are \emph{isomorphic},
written $E\cong F$, if there exists an isomorphism $f\colon E\to F$.
\end{definition}

\begin{theorem}{Classification by dimension}{thm:classification-dim}
\index{isomorphism!classification}
Two finite-dimensional $\K$-vector spaces are isomorphic if and only
if they have the same dimension. In particular, every
$n$-dimensional $\K$-vector space is isomorphic to~$\K^n$.
\end{theorem}

\begin{proof}
$(\Rightarrow)$\; If $f\colon E\isoto F$ is an isomorphism and
$(\vece{1},\ldots,\vece{n})$ is a basis of~$E$, then
$(f(\vece{1}),\ldots,f(\vece{n}))$ is a basis of~$F$
(an injective linear map sends linearly independent sets to linearly
independent sets, and a surjective map ensures they Span~$F$). Hence
$\dim F = n = \dim E$.

$(\Leftarrow)$\; Let $\dim E = \dim F = n$. Choose bases
$(\vece{1},\ldots,\vece{n})$ of~$E$ and
$(\vecf{1},\ldots,\vecf{n})$ of~$F$. By

ef{prop:prop:linear-determined-by-basis}, there is a unique linear map
$f\colon E\to F$ with $f(\vece{i}) = \vecf{i}$. Since $f$ maps
a basis to a basis, it is bijective.
\end{proof}


% ============================================================
\section{Composition of linear maps}\label{sec:composition}
% ============================================================

\begin{proposition}{Composition is linear}{prop:composition-linear}
\index{composition!of linear maps}
If $f\colon E\to F$ and $g\colon F\to G$ are linear maps, then the
composition $g\circ f\colon E\to G$ is linear.
\end{proposition}

\begin{proof}
For all $\lambda,\mu\in\K$ and $u,v\in E$:
\[
  (g\circ f)(\lambda u+\mu v)
  = g\!\bigl(f(\lambda u+\mu v)\bigr)
  = g\!\bigl(\lambda f(u)+\mu f(v)\bigr)
  = \lambda\, g(f(u))+\mu\, g(f(v))
  = \lambda\,(g\circ f)(u)+\mu\,(g\circ f)(v).\qedhere
\]
\end{proof}

\begin{proposition}{Properties of composition}{prop:composition-props}
Let $f\colon E\to F$, $g\colon F\to G$, $h\colon G\to H$ be linear.
\begin{enumerate}[label=(\roman*)]
  \item \textbf{Associativity:} $h\circ(g\circ f) = (h\circ g)\circ f$.
  \item \textbf{Identity:} $f\circ\Id_E = f = \Id_F\circ f$.
  \item If $f$ is an isomorphism, its inverse $\inv{f}\colon F\to E$
    is also linear, with $\inv{f}\circ f = \Id_E$ and
    $f\circ\inv{f} = \Id_F$.
\end{enumerate}
\end{proposition}

\begin{proof}
(i) and (ii) follow from the corresponding properties of function
composition. For (iii), let $w_1,w_2\in F$ and $\lambda\in\K$. Write
$w_i = f(v_i)$. Then $\inv{f}(\lambda w_1 + w_2) =
\inv{f}(\lambda f(v_1)+f(v_2)) = \inv{f}(f(\lambda v_1+v_2)) =
\lambda v_1+v_2 = \lambda\inv{f}(w_1)+\inv{f}(w_2)$.
\end{proof}

\begin{proposition}{Kernel and image under composition}%
{prop:ker-im-composition}
Let $f\colon E\to F$ and $g\colon F\to G$ be linear. Then:
\begin{enumerate}[label=(\roman*)]
  \item $\Ker f \subseteq \Ker(g\circ f)$.
  \item $\im(g\circ f) \subseteq \im g$.
  \item If $g$ is injective, then $\Ker(g\circ f) = \Ker f$.
  \item If $f$ is surjective, then $\im(g\circ f) = \im g$.
\end{enumerate}
\end{proposition}

\begin{proof}
(i) If $v\in\Ker f$, then $(g\circ f)(v) = g(\vzero_F) = \vzero_G$.

(ii) Every element of $\im(g\circ f)$ has the form $g(f(v))\in\im g$.

(iii) $v\in\Ker(g\circ f)$ means $g(f(v))=\vzero_G$. If $g$ is
injective, $f(v) = \vzero_F$, so $v\in\Ker f$.

(iv) If $f$ is surjective, every $w\in F$ is $f(v)$ for some $v$,
so $g(w) = g(f(v)) \in \im(g\circ f)$. Hence $\im g\subseteq\im(g\circ f)$.
\end{proof}


% ============================================================
\section{The vector space of linear maps}\label{sec:space-of-linear-maps}
% ============================================================

\begin{theorem}{$\cL(E,F)$ is a vector space}{thm:L-E-F}
\index{linear map!space of}
Let $E$ and $F$ be vector spaces over~$\K$. Define operations on the
set $\cL(E,F)$ of all linear maps from $E$ to $F$:
\begin{itemize}
  \item \textbf{Addition:} $(f+g)(v) \deff f(v)+g(v)$.
  \item \textbf{Scalar multiplication:} $(\lambda f)(v) \deff \lambda\,f(v)$.
\end{itemize}
Then $\cL(E,F)$ is a $\K$-vector space, with zero element the zero map
$\vzero\colon v\mapsto \vzero_F$.
\end{theorem}

\begin{proof}
One must verify the vector space axioms. We check that $f+g$ and
$\lambda f$ are indeed linear (so that $\cL(E,F)$ is closed under
these operations), and the remaining axioms are inherited from those
of~$F$ applied pointwise.

\emph{$f+g$ is linear:}
$(f+g)(\alpha u + \beta v)
= f(\alpha u+\beta v) + g(\alpha u+\beta v)
= \alpha f(u)+\beta f(v)+\alpha g(u)+\beta g(v)
= \alpha(f+g)(u)+\beta(f+g)(v)$.

\emph{$\lambda f$ is linear:}
$(\lambda f)(\alpha u+\beta v)
= \lambda f(\alpha u+\beta v)
= \lambda(\alpha f(u)+\beta f(v))
= \alpha(\lambda f)(u)+\beta(\lambda f)(v)$.

Associativity of addition, commutativity, etc., all follow from the
corresponding properties in~$F$.
\end{proof}

\begin{proposition}{Dimension of $\cL(E,F)$}{prop:dim-L-E-F}
If $\dim E = n$ and $\dim F = m$, then $\dim\cL(E,F) = nm$.
\end{proposition}

\begin{proof}
This follows from the isomorphism $\cL(E,F)\cong\Mat_{m,n}(\K)$
established in 
ef{sec:matrix-of-linear-map}: the space of
$m\times n$ matrices has dimension~$mn$.
\end{proof}

\begin{remark}{The endomorphism algebra}{rem:endomorphism-algebra}
\index{endomorphism!algebra}
The space $\End(E) = \cL(E,E)$ carries an additional operation:
composition. This makes $\End(E)$ a \emph{$\K$-algebra} (a vector
space equipped with an associative bilinear multiplication and a
unit element~$\Id_E$). Note that composition is not commutative
in general.
\end{remark}


% ============================================================
\section{Matrix of a linear map}\label{sec:matrix-of-linear-map}
% ============================================================

\begin{definition}{Matrix representation}{def:matrix-representation}
\index{matrix!of a linear map}
Let $E$ and $F$ be finite-dimensional with ordered bases
$\cB = (\vece{1},\ldots,\vece{n})$ and
$\cC = (\vecf{1},\ldots,\vecf{m})$ respectively.
Let $f\colon E\to F$ be linear. For each $j\in\set{1,\ldots,n}$, write
\[
  f(\vece{j}) = \sum_{i=1}^{m} a_{ij}\,\vecf{i}.
\]
The $m\times n$ matrix
\[
  \Mat_{\cC,\cB}(f) \deff (a_{ij})_{\substack{1\le i\le m\\1\le j\le n}}
  = \begin{pmatrix}
    a_{11} & a_{12} & \cdots & a_{1n}\\
    a_{21} & a_{22} & \cdots & a_{2n}\\
    \vdots & \vdots & \ddots & \vdots\\
    a_{m1} & a_{m2} & \cdots & a_{mn}
  \end{pmatrix}
\]
is called the \emph{matrix of~$f$ relative to the bases $\cB$
and~$\cC$}. The $j$-th column is the coordinate vector of $f(\vece{j})$
in basis~$\cC$.
\end{definition}

\begin{remark}{Convention}{rem:matrix-convention}
When $E = F$ and $\cB = \cC$, we simply write $\Mat_\cB(f)$.
When working in $\K^n$ with the canonical basis~$\cE$, we identify
$f$ with the matrix $A = \Mat_\cE(f)$ and write $f(x) = Ax$.
\end{remark}

\begin{example}{Matrix of a rotation}{ex:matrix-rotation}
The rotation $R_\theta\colon\R^2\to\R^2$ of

ef{ex:ex:rotation} has matrix in the canonical basis:
\[
  \Mat_\cE(R_\theta)
  = \begin{pmatrix}\cos\theta & -\sin\theta\\
    \sin\theta & \cos\theta\end{pmatrix}.
\]
The first column is $R_\theta(\vece{1}) = (\cos\theta,\sin\theta)^{\mathsf{T}}$
and the second column is $R_\theta(\vece{2}) = (-\sin\theta,\cos\theta)^{\mathsf{T}}$.
\end{example}

\begin{example}{Matrix of differentiation}{ex:matrix-diff}
Consider $D\colon\cP_3(\R)\to\cP_3(\R)$, $D(p) = p'$, with basis
$\cB = (1,t,t^2,t^3)$. Then $D(1)=0$, $D(t)=1$, $D(t^2)=2t$,
$D(t^3)=3t^2$, so
\[
  \Mat_\cB(D) = \begin{pmatrix}
    0 & 1 & 0 & 0\\
    0 & 0 & 2 & 0\\
    0 & 0 & 0 & 3\\
    0 & 0 & 0 & 0
  \end{pmatrix}.
\]
\end{example}

\begin{theorem}{The matrix representation is an isomorphism}%
{thm:matrix-iso}
The map
\[
  \Phi\colon \cL(E,F) \to \Mat_{m,n}(\K), \qquad
  f \mapsto \Mat_{\cC,\cB}(f)
\]
is an isomorphism of vector spaces. Moreover:
\begin{enumerate}[label=(\roman*)]
  \item $\Mat_{\cC,\cB}(f+g) = \Mat_{\cC,\cB}(f) + \Mat_{\cC,\cB}(g)$.
  \item $\Mat_{\cC,\cB}(\lambda f) = \lambda\,\Mat_{\cC,\cB}(f)$.
  \item If $g\colon F\to G$ is linear with $G$ having basis~$\cD$, then
    \[
      \Mat_{\cD,\cB}(g\circ f)
      = \Mat_{\cD,\cC}(g)\cdot\Mat_{\cC,\cB}(f).
    \]
\end{enumerate}
\end{theorem}

\begin{proof}
Linearity of $\Phi$ (parts (i) and (ii)) follows directly from the
definitions. Bijectivity follows from

ef{prop:prop:linear-determined-by-basis}: for every matrix
$A\in\Mat_{m,n}(\K)$, there is a unique $f\in\cL(E,F)$ with
$\Phi(f) = A$ (define $f$ by $f(\vece{j}) = \sum_i a_{ij}\vecf{i}$).

(iii) Let $A = \Mat_{\cC,\cB}(f)$ and $B = \Mat_{\cD,\cC}(g)$.
The $j$-th column of $\Mat_{\cD,\cB}(g\circ f)$ is the coordinate
vector of $(g\circ f)(\vece{j}) = g(f(\vece{j}))$ in basis~$\cD$.
Now $f(\vece{j}) = \sum_{k=1}^m a_{kj}\vecf{k}$, so
\[
  g(f(\vece{j})) = \sum_{k=1}^m a_{kj}\,g(\vecf{k})
  = \sum_{k=1}^m a_{kj}\sum_{i=1}^\ell b_{ik}\,\vec{d}_i
  = \sum_{i=1}^\ell\Bigl(\sum_{k=1}^m b_{ik}\,a_{kj}\Bigr)\vec{d}_i.
\]
The coefficient of $\vec{d}_i$ is $\sum_k b_{ik}a_{kj} = (BA)_{ij}$,
proving $\Mat_{\cD,\cB}(g\circ f) = BA$.
\end{proof}

\begin{remark}{Matrix multiplication = composition}{rem:mult-comp}
\index{matrix!multiplication!and composition}

ef{thm:thm:matrix-iso}(iii) reveals the true meaning of matrix
multiplication: it is the algebraic encoding of the composition of
linear maps. The product $BA$ of matrices is defined precisely so that
the matrix of $g\circ f$ equals the product of the matrices of $g$
and~$f$.
\end{remark}


% ============================================================
\section{Change of basis}\label{sec:change-of-basis}
% ============================================================

\begin{definition}{Transition matrix (change-of-basis matrix)}%
{def:transition-matrix}
\index{transition matrix}\index{change of basis}
Let $\cB = (\vece{1},\ldots,\vece{n})$ and
$\cB' = (\vece{1}',\ldots,\vece{n}')$ be two bases of~$E$.
The \emph{transition matrix} from $\cB$ to $\cB'$ is
\[
  P_{\cB\to\cB'} \deff \Mat_{\cB}(\Id_E,\cB',\cB)
  = \bigl[\,[\vece{1}']_\cB \;\; [\vece{2}']_\cB \;\; \cdots \;\;
    [\vece{n}']_\cB\,\bigr],
\]
where $[\vece{j}']_\cB$ is the column of coordinates of $\vece{j}'$
in basis~$\cB$.
\end{definition}

\begin{proposition}{Properties of transition matrices}%
{prop:transition-props}
\begin{enumerate}[label=(\roman*)]
  \item $P_{\cB\to\cB'}$ is invertible, with
    $\inv{(P_{\cB\to\cB'})} = P_{\cB'\to\cB}$.
  \item If $[v]_\cB$ and $[v]_{\cB'}$ denote the coordinate
    vectors of $v\in E$ in the two bases, then
    \[
      [v]_\cB = P_{\cB\to\cB'}\,[v]_{\cB'}.
    \]
  \item For three bases $\cB,\cB',\cB''$:
    $P_{\cB\to\cB''} = P_{\cB\to\cB'}\cdot P_{\cB'\to\cB''}$.
\end{enumerate}
\end{proposition}

\begin{proof}
(i) The transition matrix is the matrix of $\Id_E$ relative to
$\cB'$ (source) and $\cB$ (target). Since $\Id_E$ is an isomorphism,
its matrix is invertible. The inverse is the matrix of $\Id_E$ relative
to $\cB$ and $\cB'$, which is $P_{\cB'\to\cB}$.

(ii) Let $v = \sum_j x_j'\,\vece{j}' = \sum_i x_i\,\vece{i}$. Then
$[v]_\cB = (x_1,\ldots,x_n)^{\mathsf{T}}$ and
$[v]_{\cB'} = (x_1',\ldots,x_n')^{\mathsf{T}}$. Since
$\vece{j}' = \sum_i p_{ij}\vece{i}$ where $P = (p_{ij})$,
$v = \sum_j x_j' \sum_i p_{ij}\vece{i} = \sum_i(\sum_j p_{ij}x_j')\vece{i}$,
giving $x_i = \sum_j p_{ij}x_j'$, i.e., $[v]_\cB = P[v]_{\cB'}$.

(iii) Follows from (i) and the multiplicativity of matrix
representations under composition:
$\Id_E = \Id_E\circ\Id_E$.
\end{proof}

\begin{theorem}{Change-of-basis formula for linear maps}%
{thm:change-of-basis}
\index{change of basis!formula}
Let $f\in\End(E)$. Let $\cB$ and $\cB'$ be bases of~$E$, and let
$P = P_{\cB\to\cB'}$. If $A = \Mat_\cB(f)$ and $A' = \Mat_{\cB'}(f)$,
then
\[
  \boxed{A' = \inv{P}\, A\, P.}
\]
More generally, if $f\colon E\to F$ with bases $\cB,\cB'$ of~$E$ and
$\cC,\cC'$ of~$F$, and $P = P_{\cB\to\cB'}$, $Q = P_{\cC\to\cC'}$, then
\[
  \Mat_{\cC',\cB'}(f) = \inv{Q}\,\Mat_{\cC,\cB}(f)\,P.
\]
\end{theorem}

\begin{proof}
We have $f = \Id_F\circ f\circ\Id_E$. Using

ef{thm:thm:matrix-iso}(iii):
\begin{align*}
  \Mat_{\cC',\cB'}(f)
  &= \Mat_{\cC',\cC'}(\Id_F)\cdot\Mat_{\cC',\cB'}(f) \\
  &= \Mat_{\cC',\cC}(\Id_F)\cdot\Mat_{\cC,\cB}(f)\cdot
    \Mat_{\cB,\cB'}(\Id_E) \\
  &= \inv{Q}\cdot A \cdot P.
\end{align*}
Here we used $\Mat_{\cC',\cC}(\Id_F) = \inv{(P_{\cC\to\cC'})} = \inv{Q}$
and $\Mat_{\cB,\cB'}(\Id_E) = P_{\cB\to\cB'} = P$.

For the special case $E = F$, $\cB = \cC$, $\cB' = \cC'$, $P = Q$:
$A' = \inv{P}AP$.
\end{proof}

\begin{definition}{Similar matrices}{def:similar}
\index{similar matrices}
Two matrices $A,A'\in\Mat_n(\K)$ are \emph{similar} if there exists an
invertible matrix $P\in\GL_n(\K)$ such that $A' = \inv{P}AP$. Similar
matrices represent the same endomorphism in different bases.
\end{definition}

\begin{example}{Change of basis for a projection}{ex:change-basis-proj}
Consider the projection $\pi\colon\R^2\to\R^2$ onto the $x$-axis,
with canonical matrix $A = \begin{psmallmatrix}1&0\\0&0\end{psmallmatrix}$.
In the basis $\cB' = \bigl((1,1)^{\mathsf{T}},\,(1,-1)^{\mathsf{T}}\bigr)$,
the transition matrix is
$P = \begin{psmallmatrix}1&1\\1&-1\end{psmallmatrix}$, with
$\inv{P} = \tfrac{1}{2}\begin{psmallmatrix}1&1\\1&-1\end{psmallmatrix}$.
Then
\[
  A' = \inv{P}AP
  = \frac{1}{2}\begin{pmatrix}1&1\\1&-1\end{pmatrix}
    \begin{pmatrix}1&0\\0&0\end{pmatrix}
    \begin{pmatrix}1&1\\1&-1\end{pmatrix}
  = \frac{1}{2}\begin{pmatrix}1&1\\1&-1\end{pmatrix}
    \begin{pmatrix}1&1\\0&0\end{pmatrix}
  = \begin{pmatrix}\tfrac{1}{2}&\tfrac{1}{2}\\[3pt]
    \tfrac{1}{2}&\tfrac{1}{2}\end{pmatrix}.
\]
\end{example}


% ============================================================
\section{Operations on matrices}\label{sec:matrix-operations}
% ============================================================

We now consolidate the algebraic operations on matrices, all of which
are motivated by the corresponding operations on linear maps.

\begin{definition}{Matrix addition and scalar multiplication}%
{def:matrix-add-scalar}
\index{matrix!addition}\index{matrix!scalar multiplication}
Let $A = (a_{ij})$ and $B = (b_{ij})$ be $m\times n$ matrices over~$\K$.
\begin{enumerate}[label=(\roman*)]
  \item $A + B \deff (a_{ij}+b_{ij})$.
  \item $\lambda A \deff (\lambda\,a_{ij})$ for $\lambda\in\K$.
\end{enumerate}
With these operations, $\Mat_{m,n}(\K)$ is a vector space of
dimension~$mn$.
\end{definition}

\begin{definition}{Matrix multiplication}{def:matrix-mult}
\index{matrix!multiplication}
Let $A = (a_{ij})\in\Mat_{m,n}(\K)$ and
$B = (b_{jk})\in\Mat_{n,p}(\K)$. Their product is the matrix
$AB = (c_{ik})\in\Mat_{m,p}(\K)$ where
\[
  c_{ik} = \sum_{j=1}^{n} a_{ij}\,b_{jk}.
\]
This is the ``row-by-column'' rule: the $(i,k)$-entry of $AB$ is the
dot product of the $i$-th row of $A$ with the $k$-th column of~$B$.
\end{definition}

\begin{proposition}{Properties of matrix multiplication}%
{prop:matrix-mult-props}
\begin{enumerate}[label=(\roman*)]
  \item \textbf{Associativity:} $A(BC) = (AB)C$ (when dimensions match).
  \item \textbf{Distributivity:} $A(B+C) = AB+AC$ and $(A+B)C = AC+BC$.
  \item \textbf{Scalar compatibility:} $\lambda(AB) = (\lambda A)B = A(\lambda B)$.
  \item \textbf{Identity:} $I_m A = A = A I_n$ for
    $A\in\Mat_{m,n}(\K)$.
  \item \textbf{Non-commutativity:} In general, $AB\neq BA$.
\end{enumerate}
\end{proposition}

\begin{definition}{Transpose}{def:transpose}
\index{transpose}
The \emph{transpose} of $A = (a_{ij})\in\Mat_{m,n}(\K)$ is
$\transpose{A} = (a_{ji})\in\Mat_{n,m}(\K)$.
\end{definition}

\begin{proposition}{Properties of the transpose}{prop:transpose-props}
\begin{enumerate}[label=(\roman*)]
  \item $\transpose{(A+B)} = \transpose{A}+\transpose{B}$.
  \item $\transpose{(\lambda A)} = \lambda\,\transpose{A}$.
  \item $\transpose{(\transpose{A})} = A$.
  \item $\transpose{(AB)} = \transpose{B}\,\transpose{A}$
    \quad (reversal of order).
\end{enumerate}
\end{proposition}

\begin{proof}
We prove (iv); the others are immediate.
Let $A\in\Mat_{m,n}(\K)$, $B\in\Mat_{n,p}(\K)$. The $(k,i)$-entry of
$\transpose{(AB)}$ is $(AB)_{ik} = \sum_j a_{ij}b_{jk}$.
The $(k,i)$-entry of $\transpose{B}\,\transpose{A}$ is
$\sum_j (\transpose{B})_{kj}(\transpose{A})_{ji}
= \sum_j b_{jk}\,a_{ij} = \sum_j a_{ij}b_{jk}$.
\end{proof}


% ============================================================
\section{Inverse matrices}\label{sec:inverse-matrices}
% ============================================================

\begin{definition}{Invertible matrix}{def:invertible}
\index{matrix!invertible}\index{invertible matrix}
A square matrix $A\in\Mat_n(\K)$ is \emph{invertible} (or
\emph{non-singular}) if there exists a matrix $B\in\Mat_n(\K)$ such
that $AB = BA = I_n$. The matrix $B$ is unique and is denoted
$\inv{A}$.
\end{definition}

\begin{proposition}{Properties of inverses}{prop:inverse-props}
Let $A,B\in\GL_n(\K)$ (i.e., $A$ and $B$ are invertible).
\begin{enumerate}[label=(\roman*)]
  \item $\inv{(\inv{A})} = A$.
  \item $\inv{(AB)} = \inv{B}\,\inv{A}$ \quad (reversal of order).
  \item $\inv{(\transpose{A})} = \transpose{(\inv{A})}$.
  \item $\inv{(\lambda A)} = \lambda^{-1}\inv{A}$ for $\lambda\neq 0$.
\end{enumerate}
\end{proposition}

\begin{proof}
(ii) We verify: $(AB)(\inv{B}\inv{A}) = A(B\inv{B})\inv{A}
= AI_n\inv{A} = A\inv{A} = I_n$, and similarly
$(\inv{B}\inv{A})(AB) = I_n$.
The remaining parts are proved similarly.
\end{proof}

\begin{theorem}{Invertibility criteria}{thm:invertibility}
\index{matrix!invertibility criteria}
For $A\in\Mat_n(\K)$, the following are equivalent:
\begin{enumerate}[label=(\roman*)]
  \item $A$ is invertible.
  \item The linear map $x\mapsto Ax$ is an isomorphism of $\K^n$.
  \item $\Ker(A) = \set{\vzero}$ (i.e., $Ax = 0$ implies $x = 0$).
  \item The columns of $A$ form a basis of $\K^n$.
  \item $\rank(A) = n$.
\end{enumerate}
\end{theorem}

\begin{proof}
This follows from 
ef{prop:prop:iso-finite-dim} applied to the linear map
$f_A\colon\K^n\to\K^n$, $f_A(x) = Ax$. The matrix $A$ is invertible if
and only if $f_A$ is bijective (i.e., an isomorphism). The equivalences
(ii)--(v) were established in 
ef{prop:prop:iso-finite-dim} and the
preceding discussion.
\end{proof}

\begin{example}{Inverse of a $2\times 2$ matrix}{ex:inverse-2x2}
\index{matrix!inverse!$2\times 2$}
For $A = \begin{psmallmatrix}a&b\\c&d\end{psmallmatrix}$ with
$ad-bc\neq 0$:
\[
  \inv{A} = \frac{1}{ad-bc}\begin{pmatrix}d & -b\\-c & a\end{pmatrix}.
\]
\end{example}


% ============================================================
\section{Dual spaces}\label{sec:dual-spaces}
% ============================================================

\begin{definition}{Dual space}{def:dual-space}
\index{dual space}\index{linear form}
Let $E$ be a vector space over~$\K$. The \emph{dual space} of $E$ is
\[
  E^* \deff \cL(E,\K) = \setcond{\vphi\colon E\to\K}{\vphi
    \text{ is linear}}.
\]
Elements of $E^*$ are called \emph{linear forms} (or \emph{covectors}
or \emph{linear functionals}).
\end{definition}

\begin{example}{Linear forms on $\R^n$}{ex:linear-form-Rn}
Every linear form on $\R^n$ has the form
$\vphi(\vec{x}) = a_1 x_1 + a_2 x_2 + \cdots + a_n x_n$
for some fixed scalars $a_1,\ldots,a_n\in\R$.
\end{example}

\begin{definition}{Dual basis}{def:dual-basis}
\index{dual basis}
Let $\cB = (\vece{1},\ldots,\vece{n})$ be a basis of the
finite-dimensional space~$E$. For each $i\in\set{1,\ldots,n}$, define
the linear form $\vece{i}^*\in E^*$ by
\[
  \vece{i}^*(\vece{j}) = \delta_{ij}
  = \begin{cases}1 & \text{if } i=j,\\0 & \text{if } i\neq j.\end{cases}
\]
The family $\cB^* = (\vece{1}^*,\ldots,\vece{n}^*)$ is called the
\emph{dual basis} of~$\cB$.
\end{definition}

\begin{theorem}{The dual basis is a basis of $E^*$}{thm:dual-basis}
If $\dim E = n$, then $\cB^*$ is a basis of $E^*$, and
$\dim E^* = n$. Every linear form $\vphi\in E^*$ can be written
uniquely as
\[
  \vphi = \sum_{i=1}^{n}\vphi(\vece{i})\,\vece{i}^*.
\]
\end{theorem}

\begin{proof}
\textit{Spanning.}\; Let $\vphi\in E^*$ and set
$\psi = \sum_{i=1}^n \vphi(\vece{i})\,\vece{i}^*$. For each basis
vector $\vece{j}$:
$\psi(\vece{j}) = \sum_i \vphi(\vece{i})\,\vece{i}^*(\vece{j})
= \sum_i \vphi(\vece{i})\,\delta_{ij} = \vphi(\vece{j})$.
Since $\vphi$ and $\psi$ agree on a basis, they are equal
(
ef{prop:prop:linear-determined-by-basis}).

\textit{Linear independence.}\; Suppose
$\sum_i \lambda_i\,\vece{i}^* = 0$ (the zero form). Evaluating at
$\vece{j}$: $\sum_i \lambda_i\,\delta_{ij} = \lambda_j = 0$ for
each~$j$.
\end{proof}

\begin{remark}{Finite dimension is essential}{rem:dual-infinite}
In infinite dimension, $E^*$ is generally much larger than~$E$.
For instance, if $E = \K[t]$ (polynomials of all degrees), then
$E^*$ is not isomorphic to~$E$. The theory of dual spaces becomes
significantly richer (and subtler) in that setting.
\end{remark}


% ============================================================
\section{Applications}\label{sec:applications-ch2}
% ============================================================

\subsection{Computer graphics transformations}%
\label{subsec:computer-graphics}
\index{computer graphics}

Every 2D transformation used in computer graphics---translation,
rotation, scaling, shearing---can be expressed as a linear map (or, in
the case of translations, an \emph{affine} map made linear by
\emph{homogeneous coordinates}).

In homogeneous coordinates, a point $(x,y)$ is represented as the
triple $(x,y,1)^{\mathsf{T}}$, and transformations become
$3\times 3$ matrices:
\[
  \underbrace{\begin{pmatrix}
    \cos\theta & -\sin\theta & t_x\\
    \sin\theta & \cos\theta & t_y\\
    0 & 0 & 1
  \end{pmatrix}}_{\text{rotate by $\theta$, translate by $(t_x,t_y)$}}
  \begin{pmatrix}x\\y\\1\end{pmatrix}
  = \begin{pmatrix}
    x\cos\theta - y\sin\theta + t_x\\
    x\sin\theta + y\cos\theta + t_y\\
    1
  \end{pmatrix}.
\]

Composing transformations corresponds to multiplying their matrices.
This is the foundation of the rendering pipeline in OpenGL, DirectX,
and every modern graphics system: the entire scene geometry is
transformed by chains of matrix multiplications.

\subsection{Preview: Markov chains}\label{subsec:markov-preview}
\index{Markov chain}

A \emph{Markov chain} models a system that transitions between finitely
many states with fixed probabilities. The transition probabilities are
encoded in a \emph{stochastic matrix}
$P = (p_{ij})\in\Mat_n(\R)$ where $p_{ij}\geq 0$ and
$\sum_j p_{ij} = 1$ for each row~$i$.

If $\vec{x}^{(k)}\in\R^n$ is the probability distribution at time
step~$k$ (a row vector), then
\[
  \vec{x}^{(k+1)} = \vec{x}^{(k)}\cdot P.
\]
The long-term behavior of the chain is governed by the eigenvalues and
eigenvectors of~$P$---a topic we will study in detail in

ef{ch:eigenvalues}. For now, observe that computing the state after
$k$ steps amounts to computing $\vec{x}^{(0)}\cdot P^k$, a purely
linear-algebraic problem.


% ============================================================
\section{Exercises}\label{sec:exercises-ch2}
% ============================================================

\begin{exercise}{Verifying linearity \basic}{exo:verify-linearity}
Which of the following maps are linear? Justify each answer.
\begin{enumerate}[label=(\alph*)]
  \item $f\colon\R^2\to\R$, $f(x,y) = 3x - 2y$.
  \item $g\colon\R^2\to\R$, $g(x,y) = x^2 + y$.
  \item $h\colon\R^3\to\R^2$, $h(x,y,z) = (x+z,\, 2y-x)$.
  \item $k\colon\R^2\to\R^2$, $k(x,y) = (x+1,\, y)$.
\end{enumerate}
\end{exercise}

\begin{exercise}{Kernel and image computation \basic}{exo:ker-im-compute}
Let $f\colon\R^3\to\R^2$ be defined by $f(x,y,z) = (x+y-z,\; 2x+y+z)$.
\begin{enumerate}[label=(\alph*)]
  \item Find $\Ker f$ and give a basis.
  \item Find $\im f$ and determine $\rank(f)$.
  \item Verify the rank--nullity theorem for this map.
\end{enumerate}
\end{exercise}

\begin{exercise}{Matrix of a linear map \basic}{exo:matrix-linear-map}
Let $f\colon\R^3\to\R^2$ be defined by
$f(\vece{1}) = (1,2)^{\mathsf{T}}$,
$f(\vece{2}) = (0,-1)^{\mathsf{T}}$,
$f(\vece{3}) = (3,0)^{\mathsf{T}}$
(with respect to the canonical bases). Write down $\Mat(f)$ and
compute $f(1,1,1)^{\mathsf{T}}$.
\end{exercise}

\begin{exercise}{Composition and matrix multiplication \intermediate}%
{exo:composition-matrices}
Let $f\colon\R^2\to\R^3$ and $g\colon\R^3\to\R^2$ have matrices
\[
  A = \begin{pmatrix}1&0\\2&1\\0&3\end{pmatrix}, \qquad
  B = \begin{pmatrix}1&-1&2\\0&1&1\end{pmatrix}.
\]
\begin{enumerate}[label=(\alph*)]
  \item Compute $BA$ and $AB$.
  \item Which product represents $g\circ f$? Which represents $f\circ g$?
  \item Determine $\Ker(g\circ f)$ and $\rank(g\circ f)$.
\end{enumerate}
\end{exercise}

\begin{exercise}{Change of basis \intermediate}{exo:change-basis}
Let $f\colon\R^2\to\R^2$ be defined by $f(x,y) = (3x+y,\; x+3y)$.
\begin{enumerate}[label=(\alph*)]
  \item Write $\Mat_\cE(f)$ in the canonical basis $\cE$.
  \item Let $\cB' = \bigl((1,1)^{\mathsf{T}},(1,-1)^{\mathsf{T}}\bigr)$.
    Find the transition matrix $P$ from $\cE$ to~$\cB'$.
  \item Compute $\Mat_{\cB'}(f) = \inv{P}\Mat_\cE(f)\,P$.
  \item Interpret the result geometrically.
\end{enumerate}
\end{exercise}

\begin{exercise}{Rank--nullity applications \intermediate}%
{exo:rank-nullity-app}
\begin{enumerate}[label=(\alph*)]
  \item Let $f\colon\R^5\to\R^3$ be linear with $\dim(\Ker f) = 2$.
    What is $\rank(f)$?
  \item Let $g\colon\R^4\to\R^4$ be linear with $\rank(g) = 3$.
    What is $\dim(\Ker g)$? Is $g$ surjective?
  \item Prove that there is no surjective linear map from $\R^3$ to $\R^4$.
  \item Prove that there is no injective linear map from $\R^4$ to $\R^3$.
\end{enumerate}
\end{exercise}

\begin{exercise}{Projections \intermediate}{exo:projections}
\index{projection!idempotent}
A linear map $p\colon E\to E$ is called a \emph{projection} if
$p\circ p = p$ (idempotent).
\begin{enumerate}[label=(\alph*)]
  \item Show that $\im p = \Ker(\Id_E - p)$.
  \item Show that $\Ker p = \im(\Id_E - p)$.
  \item Deduce that $E = \Ker p \oplus \im p$.
  \item If $\dim E = n$ and $\rank(p) = r$, show there exists a basis
    in which $\Mat(p) = \diag(1,\ldots,1,0,\ldots,0)$ ($r$ ones).
\end{enumerate}
\end{exercise}

\begin{exercise}{Dual basis computation \intermediate}{exo:dual-basis}
Let $\cB = (\vece{1},\vece{2})$ be the basis of $\R^2$ with
$\vece{1} = (1,1)^{\mathsf{T}}$, $\vece{2} = (1,-1)^{\mathsf{T}}$.
Find explicit formulas for $\vece{1}^*$ and $\vece{2}^*$ (as functions
of the canonical coordinates $x_1,x_2$).
\end{exercise}

\begin{exercise}{Injectivity and surjectivity \challenging}%
{exo:inject-surject}
Let $f\colon E\to F$ and $g\colon F\to G$ be linear maps between
finite-dimensional spaces.
\begin{enumerate}[label=(\alph*)]
  \item Prove: $\rank(g\circ f) \leq \min(\rank f,\, \rank g)$.
  \item Prove the \emph{Sylvester rank inequality}\index{Sylvester rank inequality}:
    $\rank f + \rank g - \dim F \leq \rank(g\circ f)$.
    \textit{Hint:} Apply the rank--nullity theorem to the restriction
    $g\restrict{\im f}$.
  \item Deduce: if $g\circ f = 0$, then
    $\rank f + \rank g \leq \dim F$.
\end{enumerate}
\end{exercise}

\begin{exercise}{Nilpotent endomorphisms \challenging}{exo:nilpotent}
\index{nilpotent}
An endomorphism $f\in\End(E)$ is \emph{nilpotent} of index~$k$ if
$f^k = 0$ but $f^{k-1}\neq 0$.
\begin{enumerate}[label=(\alph*)]
  \item Show that if $f$ is nilpotent, then $f$ is not invertible.
  \item Show that the chain $\set{\vzero}\subseteq\Ker f\subseteq
    \Ker f^2\subseteq\cdots$ is strictly increasing until it
    stabilizes at~$E$.
  \item Prove that the nilpotency index satisfies $k\leq\dim E$.
  \item Prove that $\Id + f$ is invertible, and find a formula for
    $(\Id + f)^{-1}$ in terms of powers of $f$.
    \textit{Hint:} geometric series.
\end{enumerate}
\end{exercise}

\begin{exercise}{Linear maps on polynomial spaces \challenging}%
{exo:polynomial-maps}
Let $E = \cP_3(\R)$ (polynomials of degree $\leq 3$) with basis
$\cB = (1,t,t^2,t^3)$. Define $T\colon E\to E$ by $T(p)(t) = p(t+1)$.
\begin{enumerate}[label=(\alph*)]
  \item Verify that $T$ is linear.
  \item Compute the matrix $\Mat_\cB(T)$.
  \item Determine $\Ker T$ and $\im T$. Is $T$ an isomorphism?
  \item Compute $T^{-1}$ and interpret it.
\end{enumerate}
\end{exercise}

\begin{exercise}{The trace as a linear form \basic}{exo:trace-linear}
\index{trace}
Show that the trace map $\tr\colon\Mat_n(\K)\to\K$, defined by
$\tr(A) = \sum_{i=1}^n a_{ii}$, is a linear form.
Determine $\Ker(\tr)$ and its dimension.
\end{exercise}


% ============================================================
\section*{Chapter summary}\label{sec:summary-ch2}
\addcontentsline{toc}{section}{Chapter summary}
% ============================================================

\begin{itemize}
  \item A \textbf{linear map}\index{linear map} $f\colon E\to F$
    preserves addition and scalar multiplication. It is completely
    determined by its values on a basis
    (
ef{prop:prop:linear-determined-by-basis}).

  \item The \textbf{kernel} $\Ker f$ and \textbf{image} $\im f$ are
    subspaces (
ef{thm:thm:ker-im-subspaces}). The kernel measures
    ``information lost'' by $f$; the image measures ``information
    reached.''

  \item The \textbf{rank--nullity theorem}
    (
ef{thm:thm:rank-nullity}):
    $\dim E = \dim(\Ker f) + \rank(f)$.

  \item A linear map between spaces of equal finite dimension is
    injective if and only if it is surjective if and only if it is
    bijective (
ef{prop:prop:iso-finite-dim}).

  \item The set $\cL(E,F)$ of all linear maps is itself a vector space
    of dimension $\dim E\cdot\dim F$
    (
ef{thm:thm:L-E-F,prop:prop:dim-L-E-F}).

  \item Choosing bases allows us to represent every linear map as a
    \textbf{matrix} (
ef{def:def:matrix-representation}).
    \textbf{Composition} of linear maps corresponds to
    \textbf{matrix multiplication}
    (
ef{thm:thm:matrix-iso}).

  \item Under a \textbf{change of basis} with transition matrix~$P$,
    the matrix of an endomorphism transforms as
    $A' = \inv{P}AP$ (
ef{thm:thm:change-of-basis}).

  \item A square matrix is \textbf{invertible} if and only if its
    kernel is trivial, equivalently if and only if its rank equals
    its size (
ef{thm:thm:invertibility}).

  \item The \textbf{dual space} $E^* = \cL(E,\K)$ has the same
    dimension as~$E$ in finite dimension, and admits a canonical
    \textbf{dual basis} (
ef{thm:thm:dual-basis}).
\end{itemize}
